\documentclass[10pt,a4paper]{amsart}
\usepackage[margin=19mm]{geometry}
\usepackage{amsmath,amssymb,mathtools}
\usepackage[scr=rsfs,cal=euler]{mathalfa}
\usepackage{xfrac}
\usepackage[unicode,hidelinks,bookmarks=false]{hyperref}
\newcommand{\ie}{i.e.\ }
\newcommand{\cf}{cf.\ }
\newcommand{\resp}{respectively\ }
\newcommand{\Subsection}{\subsection}
\providecommand{\nobreakdash}{\nobreak\hbox{-}}
\setlength{\emergencystretch}{2em}
\multlinegap0pt
\usepackage{amssymb}
\usepackage{algorithm}\usepackage{algpseudocode}
\usepackage{mathtools}
\usepackage[shortlabels]{enumitem}
\usepackage{bm}
\usepackage{xcolor}
\usepackage[normalem]{ulem}
\newtheorem{lemma}{Lemma}
\newtheorem{corollary}{Corollary}
\newtheorem{theorem}{Theorem}
\usepackage{cleveref}
\crefname{lemma}{Lemma}{Lemmas}
\crefname{corollary}{Corollary}{Corollarys}
\crefname{theorem}{Theorem}{Theorems}
\newcommand{\ts}{^T}
\title{A Variational Equation and Lower Bound for the Linear Least-Squares Backward Error}
\author{Eric Hallman}
\date{Source: arXiv:2605.09211v1 (CC BY 4.0)}
\begin{document}
\maketitle

\noindent\emph{Source and licence.} A Variational Equation and Lower Bound for the Linear Least-Squares Backward Error. Version arXiv:2605.09211v1 (CC BY 4.0), distributed by arXiv under the Creative Commons Attribution 4.0 International licence, http://creativecommons.org/licenses/by/4.0/, as recorded in the arXiv licence metadata for this exact version and shown on the abstract page https://arxiv.org/abs/2605.09211v1. Full licence notice: \texttt{licenses/arxiv-2605.09211-CC-BY-4.0.txt}.\par\smallskip

\noindent\textit{Editorial scope.} Counted unit: the article's theorem thm:rank\_two\_eigenvalue\_problem (source0.tex lines 422--428). The article's complete original proof of it is delivered with it (source0.tex lines 429--464, one \texttt{proof} environment of 36 lines). No mathematics is written, repaired or re-derived: the statement and the proof are the article's own lines, in the article's own notation, and no retained line is substituted or re-flowed.  Retained uncounted as the source-local context the proof needs: lines 160--162; lines 202--211; lines 308--313; lines 323--325; lines 395--400.  Nothing was omitted from any retained span, so the package carries no omission record.  No retained line contains a citation, so the wrapper carries no bibliography and no bibliography entry.  No retained line contains a figure environment, an image-inclusion command or drawing code, so no image is delivered and the package declares no asset dependency.  Editorial formatting only, in addition: an AMS \texttt{amsart} preamble that restates the article's own declarations and macros used by the retained lines, namely the environments \texttt{lemma}, \texttt{corollary}, \texttt{theorem} on separate counters, together with the standard packages those lines need.  The retained environments print in this document as Lemma 1, Corollary 1, Theorem 1 in order of appearance, whereas the article prints the same items with the numbering of its own full document.  The complete native source of the article as distributed in the arXiv e-print for this exact version is bundled unchanged as \texttt{sources/200208/source0.tex}.
\begin{equation}\label{def:mu_min_proj}
    \mu = \mu(A,R_\theta) \equiv \min_{Y} \left(\|YY^\dagger A\|_F^2 + \|(I-YY^\dagger)R_\theta\|_F^2\right)^{\frac{1}{2}}. 
\end{equation}

\begin{lemma}\label{lemma:continuity_argument}
    The backward error as defined in \cref{def:mu_min_proj} satisfies
    \[
    \lim_{\epsilon \rightarrow 0}\mu\left(\begin{bmatrix}
        A \\ \epsilon I \\ 0
    \end{bmatrix}, \begin{bmatrix}
        R_\theta \\ 0 \\ \epsilon I
    \end{bmatrix}\right) = \mu(A, R_\theta). 
    \]
\end{lemma}

\begin{lemma}\label{lemma:variational_first}
For any $A\in \mathbb{R}^{m\times n}$ and $R_\theta\in \mathbb{R}^{m\times d}$,
\begin{equation}\label{eqn:variational_first}
    \mu(A, R_\theta) = \sup_{X\ts J_{n,d}X = -I_d} \left(\|R_\theta\|_F^2 - \|[A, R_\theta]X\|_F^2\right)^{\frac{1}{2}}.
\end{equation}
\end{lemma}

    \begin{equation}\label{eqn:negative_genl_eigvals}
        \min_{X\ts J_{n,d}X = -I_d} \|[A,R_\theta]X\|_F^2 = -\sum_{i=1}^d\lambda_i^{-},
    \end{equation}

\begin{corollary}\label{cor:backward_error_fixed_point}
    For $A\in \mathbb{R}^{m\times n}$ and $r_\theta\in \mathbb{R}^m$, $\mu(A,r_\theta)$ is the smallest nonnegative number solving
    \begin{equation*}
        \mu^2 = r_\theta\ts A(A\ts A + (\|r_\theta\|_2^2 - \mu^2)I)^{-1}A\ts r_\theta.
    \end{equation*}
\end{corollary}

\begin{theorem}\label{thm:rank_two_eigenvalue_problem}
    For $a\in \mathbb{R}^m$ and $r \in \mathbb{R}^m$, at least one of which is nonzero, 
    \begin{equation*}
        \mu(a,r) = \frac{2 |a\ts r|}{\|a+r\|_2 + \|a-r\|_2},
    \end{equation*}
    where $\mu(a,r)$ is as defined in \cref{def:mu_min_proj}.
\end{theorem}

\begin{proof}
    By \Cref{cor:backward_error_fixed_point}, $\mu(a,r)$ is the smallest nonnegative solution to the equation
    \begin{equation*}
        \mu^2 = \frac{(a\ts r)^2}{\|a\|_2^2 + \|r\|_2^2 - \mu^2}.
    \end{equation*}
    % Assume that $a$ and $r$ are linearly independent; the linearly dependent case will follow by a continuity argument as in \cref{lemma:continuity_argument}. By \cref{lemma:variational_first} and \cref{eqn:negative_genl_eigvals}, we know that
    % \begin{equation} \label{eqn:mu_shifted_eigenvalue}
    %     \mu(a,r) = (\|r\|_2^2 + \lambda_1^-)^{\frac{1}{2}}, 
    % \end{equation}
    % where $\lambda_1^- < 0$, being a generalized eigenvalue of $([a,r]\ts [a,r], J_{1,1})$, satisfies
    % \begin{equation*}
    %     \det\left(\begin{bmatrix}
    %         \|a\|_2^2 & a\ts r\\
    %         r\ts a & \|r\|_2^2
    %     \end{bmatrix} - \lambda_1^{-}J_{1,1}\right) = 0.
    % \end{equation*}
    Solving the resulting quadratic equation yields 
    % \begin{align*}
    %     \lambda_1^- &= \frac{\|a\|_2^2 - \|r\|_2^2 - \sqrt{(\|r\|_2^2-\|a\|_2^2)^2 + 4\|a\|_2^2\|r\|_2^2 - 4(a\ts r)^2}}{2} \\
    %     &= \frac{\|a\|_2^2 - \|r\|_2^2 - \sqrt{(\|r\|_2^2+\|a\|_2^2)^2  - 4(a\ts r)^2}}{2} \\
    %     &= \frac{\|a\|_2^2 - \|r\|_2^2 - \sqrt{\|a+r\|_2^2\|a-r\|_2^2}}{2} \\
    %     &= \frac{\|a\|_2^2 - \|r\|_2^2 - \|a+r\|_2\|a-r\|_2}{2}.
    % \end{align*}
    % Therefore, from \cref{eqn:mu_shifted_eigenvalue} we find that
    \begin{align*}
        \mu(a,r) &= \left(\frac{\|a\|_2^2 + \|r\|_2^2 -\sqrt{(\|a\|_2^2 + \|r\|_2^2)^2 - 4(a\ts r)^2}}{2}\right)^{\frac{1}{2}} \\
        &= \left(\frac{\|a\|_2^2 + \|r\|_2^2 - \|a+r\|_2\|a-r\|_2}{2}\right)^{\frac{1}{2}} \\
        &= \left(\frac{\|a+r\|_2^2 + \|a-r\|_2^2 - 2\|a+r\|_2\|a-r\|_2}{4}\right)^{\frac{1}{2}}\\
        &= \frac{|\|a+r\|_2 - \|a-r\|_2|}{2}.
    \end{align*}
    Although elegant, this formula is unstable. The stable formulation is
    \begin{equation*}
        \mu = \frac{|\|a+r\|_2 - \|a-r\|_2|}{2}\cdot \frac{\|a+r\|_2 + \|a-r\|_2}{\|a+r\|_2+\|a-r\|_2} = \frac{2|a\ts r|}{\|a+r\|_2+\|a-r\|_2},
    \end{equation*}
    which completes the proof.
\end{proof}

\end{document}
